
%%%%%%%%%%%%%%%%%%%%%%%%%%%%%%%%%%%%%%%%
%           main body
%%%%%%%%%%%%%%%%%%%%%%%%%%%%%%%%%%%%%%%%

%-----------------------------------
%	TITLE PAGE
%-----------------------------------

\title[第十一章\\
Euclid空间上的极限与连续]{\texorpdfstring{\LARGE \bfseries 第十一章\quad Euclid 空间上的极限与连续}{第十一章 Euclid 空间上的极限与连续}} % The short title appears at the bottom of every slide, the full title is only on the title page
\author[\S 11.3\\
多元连续函数性质] % Your institution as it will appear on the bottom of every slide, maybe shorthand to save space
{\Large \bfseries \S 11.3 多元连续函数的性质}
% \author{Singular Value Decomposition} % Your name
% \date{\today} % Date, can be changed to a custom date
\date{}

%------------------------------------------------

\begin{document}

%------------------------------------------------

\begin{frame}

\titlepage % Print the title page as the first slide

\end{frame}

%------------------------------------------------

% \begin{frame}
% \frametitle{内容提要} % Table of contents slide, comment this block out to remove it
% \tableofcontents % Throughout your presentation, if you choose to use \section{} and \subsection{} commands, these will automatically be printed on this slide as an overview of your presentation
% \end{frame}

%-----------------------------------
%	PRESENTATION SLIDES
%-----------------------------------

%------------------------------------------------

\begin{frame}{回忆: 闭区间上一元连续函数性质}

我们回忆一下定义在闭区间 (紧集) $[a, b]$ 上一元连续函数 $f(x)$ 的主要性质:

\pause

\begin{enumerate}
\item 
\end{enumerate}

\end{frame}

%------------------------------------------------

\section{紧集上的连续函数}

%------------------------------------------------

\begin{frame}{}



\end{frame}

%------------------------------------------------

\section{连通集}

%------------------------------------------------

\begin{frame}{连通集}



\end{frame}

%------------------------------------------------

\begin{frame}{一般拓扑空间内的连通集}



\end{frame}

%------------------------------------------------

\section{连通集上的连续映射}

%------------------------------------------------

\begin{frame}{xxx}



\end{frame}


%------------------------------------------------

\end{document}
